\section{Three rows}\label{sec:three}

On three rows the row bound of Theorem~\ref{thm:rowbound} is $m+2$, and it is attained. We
prove this first. Then we write every maximum blackout as a word in seven letters and count
the words. The counts $a_3(m)$ form row $3$ of \seqnum{A399580}, which begins $53$, $224$,
$892$, $3420$, $12704$ at widths $3$ to $7$. Theorem~\ref{thm:threecount} gives an exact
formula, Theorem~\ref{thm:threegf} the rational generating function,
Corollary~\ref{cor:threeasymp} the asymptotic $a_3(m)\sim\frac4{81}m^23^m$, and
Theorem~\ref{thm:threeclosed} a closed form of the remainder. Throughout this section
$m\ge3$, the rows are $1,2,3$ from top to bottom, and the tilted rectangles are the diamonds
$D(i,1)=\{(i,2),(i+1,1),(i+2,2),(i+1,3)\}$ for $1\le i\le m-2$
(Lemma~\ref{lem:tilted-narrow}).

\subsection{The value}

\begin{theorem}[three rows]\label{thm:threevalue}
For every $m\ge3$, the maximum blackout of the $3\times m$ grid has $m+2$ points, that is,
$f(3,m)=m+2$.
\end{theorem}

\begin{proof}
The upper bound is the case $h=3$ of Theorem~\ref{thm:rowbound}, since
$q_3=\lfloor9/4\rfloor=2$. The blackout of row $1$ and column $1$,
\[
S=\{(i,1):1\le i\le m\}\cup\{(1,2),(1,3)\},
\]
has $m+2$ points, and it is the boundary cross of Lemma~\ref{lem:cross} for $h=3$, so it is
valid.
\end{proof}

Figure~\ref{fig:three} shows $S$ for $m=6$. Its graph $G_S$ is a tree, so no two axis-aligned
rectangles agree (Lemma~\ref{lem:girth}). The diamond $D(i,1)$ shows three corners for
$i\ge2$, and $D(1,1)$ shows exactly $(3,2)$ and $(2,3)$, which no axis-aligned rectangle shows.
The trureturing project~\cite{trureturing2026} also proves the value, independently and with
a machine-checked proof.

\begin{figure}[t]
\centering
\begin{tikzpicture}[x=0.665cm,y=0.665cm,baseline=(current bounding box.north)]
  \draw[gridcol] (1,-1) grid (6,-3);
  \fill[black] (1,-1) circle (0.17);
  \fill[black] (2,-1) circle (0.17);
  \fill[black] (3,-1) circle (0.17);
  \fill[black] (4,-1) circle (0.17);
  \fill[black] (5,-1) circle (0.17);
  \fill[black] (6,-1) circle (0.17);
  \fill[black] (1,-2) circle (0.17);
  \filldraw[fill=white,draw=keptline,line width=0.6pt] (2,-2) circle (0.17);
  \filldraw[fill=white,draw=keptline,line width=0.6pt] (3,-2) circle (0.17);
  \filldraw[fill=white,draw=keptline,line width=0.6pt] (4,-2) circle (0.17);
  \filldraw[fill=white,draw=keptline,line width=0.6pt] (5,-2) circle (0.17);
  \filldraw[fill=white,draw=keptline,line width=0.6pt] (6,-2) circle (0.17);
  \fill[black] (1,-3) circle (0.17);
  \filldraw[fill=white,draw=keptline,line width=0.6pt] (2,-3) circle (0.17);
  \filldraw[fill=white,draw=keptline,line width=0.6pt] (3,-3) circle (0.17);
  \filldraw[fill=white,draw=keptline,line width=0.6pt] (4,-3) circle (0.17);
  \filldraw[fill=white,draw=keptline,line width=0.6pt] (5,-3) circle (0.17);
  \filldraw[fill=white,draw=keptline,line width=0.6pt] (6,-3) circle (0.17);
  \node[lab] at (1,-0.35) {1};
  \node[lab] at (2,-0.35) {2};
  \node[lab] at (3,-0.35) {3};
  \node[lab] at (4,-0.35) {4};
  \node[lab] at (5,-0.35) {5};
  \node[lab] at (6,-0.35) {6};
  \node[lab] at (0.35,-1) {1};
  \node[lab] at (0.35,-2) {2};
  \node[lab] at (0.35,-3) {3};
  \draw[draw=rectA,line width=1.4pt,line join=round] (2,-1) -- (3,-2) -- (2,-3) -- (1,-2) -- cycle;
  \draw[draw=rectB,line width=1.4pt,line join=round,dash pattern=on 3pt off 2pt] (4,-1) -- (5,-2) -- (4,-3) -- (3,-2) -- cycle;
  \node[capt,anchor=north,text width=4.52cm,align=center] at (3.5,-4.6) {(a) row $1$ and column $1$, with $D(1,1)$ and $D(3,1)$};
  \draw[draw=black!75,line width=0.9pt] (11.4,-0.6) -- (8.4,-2.1);
  \draw[draw=black!75,line width=0.9pt] (7.7,-3.6) -- (8.4,-2.1);
  \draw[draw=black!75,line width=0.9pt] (9.1,-3.6) -- (8.4,-2.1);
  \draw[draw=black!75,line width=0.9pt] (11.4,-0.6) -- (9.6,-2.1);
  \draw[draw=black!75,line width=0.9pt] (11.4,-0.6) -- (10.8,-2.1);
  \draw[draw=black!75,line width=0.9pt] (11.4,-0.6) -- (12.0,-2.1);
  \draw[draw=black!75,line width=0.9pt] (11.4,-0.6) -- (13.2,-2.1);
  \draw[draw=black!75,line width=0.9pt] (11.4,-0.6) -- (14.4,-2.1);
  \node[rowv] at (11.4,-0.6) {$r_{1}$};
  \node[colv] at (8.4,-2.1) {$c_{1}$};
  \node[colv] at (9.6,-2.1) {$c_{2}$};
  \node[colv] at (10.8,-2.1) {$c_{3}$};
  \node[colv] at (12.0,-2.1) {$c_{4}$};
  \node[colv] at (13.2,-2.1) {$c_{5}$};
  \node[colv] at (14.4,-2.1) {$c_{6}$};
  \node[rowv] at (7.7,-3.6) {$r_{2}$};
  \node[rowv] at (9.1,-3.6) {$r_{3}$};
  \node[capt,anchor=north,text width=4.32cm,align=center] at (11.4,-4.6) {(b) $G_S$ is a tree};
\end{tikzpicture}

\caption{The three-row construction for $m=6$. (a) Row $1$ and column $1$ are black. The diamond
$D(1,1)$ (solid) shows only $(3,2)$ and $(2,3)$, and the diamond $D(3,1)$ (dashed) shows three
corners. (b) The row-column graph $G_S$ is a tree.}
\label{fig:three}
\end{figure}

\subsection{Words and the two types}

From now on every blackout has $m+2$ points, and we describe it column by column. For a
column $i$ let $B_i\subseteq\{1,2,3\}$ be its set of black rows. We write the seven nonempty
sets as letters,
\begin{align*}
a&=\{1\},& b&=\{2\},& c&=\{3\},\\
A&=\{1,2\},& B&=\{1,3\},& C&=\{2,3\},& X&=\{1,2,3\}.
\end{align*}
So $a,b,c$ are the singletons, $A,B,C$ are the doubles, $X$ is the triple, and $b$ is the
middle row. A blackout with no empty column is the same thing as its \emph{word}
$B_1B_2\cdots B_m$. The \emph{pair set} $F$ of a word is the set of row pairs $\{r,s\}$ with
$\{r,s\}\subseteq B_i$ for some column $i$, and we name a pair by its double letter. So a
word with one $X$ has $F=\{A,B,C\}$, and a word with the doubles $A$ and $C$ has $F=\{A,C\}$.

\begin{proposition}[the two types]\label{prop:threeclass}
Let $m\ge3$ and let $S$ be a blackout of the $3\times m$ grid with $m+2$ points.
\begin{enumerate}[(a)]
\item If $S$ is valid, then every column is nonempty and the word of $S$ has one of two forms.
Either it has exactly one letter $X$ and singletons elsewhere (type T), or it has exactly two
double letters, which are different, and singletons elsewhere (type P).
\item Conversely, for every word of type T or type P, the graph $G_S$ is a tree, so no two
axis-aligned rectangles have the same presentation.
\item There are exactly $m^23^{m-1}$ words of type T or type P.
\end{enumerate}
\end{proposition}

\begin{proof}
(a) By Lemma~\ref{lem:H}, $|S|=t+e(H)$ with $t\le m$ the number of nonempty columns and $H$
the row graph, which is triangle-free on three vertices and has at most two edges
(Lemma~\ref{lem:mantel}). Since $|S|=m+2$, both bounds are attained, so $t=m$ and
$e(H)=\sum_i(|B_i|-1)=2$. So one column is a triple and the rest singletons, or two columns
are doubles and the rest singletons. Two equal doubles form a cycle of length four in $G_S$,
which Lemma~\ref{lem:girth} forbids.

(b) For type T the three row vertices hang on the column of the $X$. For type P the two
different doubles form a path through the three row vertices. In both types every other
column hangs on one row. So $G_S$ is connected with $m+3$ vertices and $m+2$ edges, a tree,
and Lemma~\ref{lem:girth} gives the claim.

(c) Type T is a choice of the column of $X$ and of $m-1$ singletons, which gives $m3^{m-1}$
words. Type P is a choice of two columns, an ordered pair of different doubles, and $m-2$
singletons, which gives $\binom m2\cdot6\cdot3^{m-2}=m(m-1)3^{m-1}$ words.
\end{proof}

We call a word of type T or type P a \emph{candidate}. So the maximum blackouts are exactly
the candidates that also keep the diamonds apart from everything else. The candidates fall
into \emph{classes} by their double letters, namely the three unordered classes $\{A,C\}$,
$\{A,B\}$ and $\{B,C\}$ of type P and the triple class $X$. The left-to-right order of the two
doubles splits a class of type P into two \emph{ordered classes}, written $AC$, $CA$ and so
on. We write $|AC_m|$ for the number of valid words of width $m$ in the ordered class $AC$,
and likewise for the other ordered classes and for $|X_m|$.

\subsection{The diamond test}

The test for a diamond is Lemma~\ref{lem:rival} read on three rows. On a candidate the answer
depends only on the three columns of the diamond and on the global pair set $F$.

\begin{lemma}[diamond test]\label{lem:diamond}
Let $m\ge3$ and let $S$ be a candidate with pair set $F$. Then $S$ is valid if and only if every
diamond $D(i,1)$, $1\le i\le m-2$, passes the test below, where $P=D(i,1)\setminus S$ is its
presentation.
\begin{enumerate}[(a)]
\item If $|P|\ge3$, the diamond passes.
\item If $P$ is empty, the diamond passes.
\item If $P=\{(i+1,1),(i+1,3)\}$, the diamond fails if and only if $B\in F$.
\item If $P=\{(i,2),(i+2,2)\}$, the diamond fails if and only if $B_i$ and $B_{i+2}$ share a
row other than $2$.
\item If $P=\{(x,r),(z,s)\}$ with $x\ne z$ and $r\ne s$, the diamond fails if and only if
$(x,s)$ and $(z,r)$ are both black.
\item If $P=\{(x,r)\}$, the diamond fails if and only if there is a row $s\ne r$ with $s\in B_x$
and $\{r,s\}\in F$.
\end{enumerate}
\end{lemma}

\begin{proof}
By Proposition~\ref{prop:threeclass}(b) no two axis-aligned rectangles agree. So $S$ is valid
if and only if no diamond has the presentation of an axis-aligned rectangle and no two
diamonds agree. A presentation with at least three points belongs to one rectangle only
(Lemma~\ref{lem:threecorners}), which is (a).

Rows (c) to (f) are the rows of Lemma~\ref{lem:rival} for a diamond. In (c) a rival lies on
rows $1,3$ and a column $y\ne i+1$ black in both rows, so it exists exactly when $B\in F$. In
(d) a rival uses columns $i,i+2$ and a row $s\ne2$ black in both columns. In (e) the two
points are opposite corners of exactly one axis-aligned rectangle, whose other corners are
$(x,s)$ and $(z,r)$. In (f) a rival uses column $x$, a row $s\ne r$ with $(x,s)$ black, and a
second column black in rows $r$ and $s$. That column is not $x$, so its existence is
$\{r,s\}\in F$. A fully black axis-aligned rectangle is a cycle of length four in $G_S$,
which a tree does not have, so an empty diamond has no axis-aligned rival. Finally two
diamonds $D(i,1)$ and $D(j,1)$ with $i<j$ share at most the tip $(i+2,2)$, when $j=i+2$. So
a common presentation has at most one point, and in both cases columns $i+1$ and $j+1$ are
black in rows $1$ and $3$, which a candidate does not allow. This proves (b).
\end{proof}

Rows (c) and (f) use the pair set of the whole word, so a double far from the diamond can
make it fail. Row (b) allows one fully black diamond, as in the valid word $bXb$ at $m=3$.

\subsection{Local rules}

A \emph{window} of a word is a string of three consecutive letters, at columns $i$, $i+1$,
$i+2$, and it \emph{passes} or \emph{fails} according to the diamond $D(i,1)$. A tip is black
exactly when its letter contains row $2$, that is, when it is $b$, $A$, $C$ or $X$.

\begin{lemma}[singleton dichotomy]\label{lem:dichotomy}
Let $S$ be a candidate and let $stu$ be a window of three singletons.
\begin{enumerate}[(a)]
\item If $B\notin F$, the window passes.
\item If $B\in F$, the window fails if and only if it is $bab$, $bbb$ or $bcb$.
\end{enumerate}
\end{lemma}

\begin{proof}
Both tips are black exactly when $s=u=b$. Then for $t=b$ the diamond shows the vertical pair
of column $i+1$, which fails exactly when $B\in F$ by row (c). For $t$ outer it shows only
the white point of column $i+1$, and by row (f) it fails exactly when $B\in F$. Suppose
exactly one tip is black, say $s=b$ and $u\ne b$. For $t=b$ three corners show. For $t$ outer
the diamond shows an oblique pair whose test in row (e) needs $(i+1,2)$ black, and it is
white. If no tip is black, three corners show.
\end{proof}

\begin{lemma}[six of nine in the class $\{A,C\}$]\label{lem:sixofnine}
Let $S$ be a candidate of class $\{A,C\}$, so $F=\{A,C\}$.
\begin{enumerate}[(a)]
\item A window $sAt$ with singletons $s,t$ fails if and only if $st$ is $bb$, $bc$ or $cb$.
\item A window $sCt$ with singletons $s,t$ fails if and only if $st$ is $bb$, $ba$ or $ab$.
\item A window with exactly one double letter, at a tip, passes.
\end{enumerate}
\end{lemma}

\begin{proof}
(a) The center $A$ makes $(i+1,1)$ black and $(i+1,3)$ white. If $s=t=b$, both tips are black
and the diamond shows only $(i+1,3)$. Row (f) with $x=i+1$, $r=3$ and the row $2\in B_{i+1}$
applies, since $\{2,3\}=C\in F$, so the window fails. If $s=b$ and $t\ne b$, the diamond shows
the oblique pair $(i+2,2)$, $(i+1,3)$, and row (e) asks whether $(i+2,3)$ and $(i+1,2)$ are
black. The second is, and the first is exactly when $t=c$. So $bAc$ fails and $bAa$ passes,
and by the mirror image $cAb$ fails and $aAb$ passes. If neither neighbor is $b$, three
corners show.

(b) The reflection in the middle row swaps $a$ with $c$ and $A$ with $C$, keeps the class
$\{A,C\}$, and maps rectangles and presentations to rectangles and presentations. So (a)
gives (b).

(c) Let the window be $Ast$, so the tip $(i,2)$ is black. If $t=b$, both tips are black. For
$s=b$ the diamond shows the vertical pair, which passes by row (c) since $B\notin F$. For $s$
outer it shows only the white point of column $i+1$, and row (f) asks whether $B\in F$. If
$t\ne b$, then for $s=b$ three corners show, and for $s$ outer the diamond shows an oblique
pair whose test in row (e) needs $(i+1,2)$ black, and it is white. The windows $Cst$, $tsA$
and $tsC$ follow by the reflection of (b) and by reversing the columns.
\end{proof}

So in the class $\{A,C\}$ an $A$ with singleton neighbors passes six of the nine possible
neighbor assignments, and Figure~\ref{fig:sixofnine} shows all nine. The rule is specific to
this class. When $B\in F$, a $B$ at the center passes only five of nine, and a double at a tip
can fail (Lemma~\ref{lem:threewindows} in Appendix~\ref{app:threegf}).

\begin{figure}[t]
\centering
\resizebox{\textwidth}{!}{\begin{tikzpicture}[x=0.52cm,y=0.52cm,baseline=(current bounding box.north)]
  \draw[gridcol] (1,-1) grid (3,-3);
  \fill[black] (1,-1) circle (0.17);
  \filldraw[fill=white,draw=keptline,line width=0.6pt] (1,-2) circle (0.17);
  \filldraw[fill=white,draw=keptline,line width=0.6pt] (1,-3) circle (0.17);
  \fill[black] (2,-1) circle (0.17);
  \fill[black] (2,-2) circle (0.17);
  \filldraw[fill=white,draw=keptline,line width=0.6pt] (2,-3) circle (0.17);
  \fill[black] (3,-1) circle (0.17);
  \filldraw[fill=white,draw=keptline,line width=0.6pt] (3,-2) circle (0.17);
  \filldraw[fill=white,draw=keptline,line width=0.6pt] (3,-3) circle (0.17);
  \draw[draw=rectA,line width=1.2pt,line join=round] (2,-1) -- (3,-2) -- (2,-3) -- (1,-2) -- cycle;
  \node[capt,anchor=north] at (2,-3.4) {$aAa$};
  \node[lab] at (0.35,-1) {1};
  \node[lab] at (0.35,-2) {2};
  \node[lab] at (0.35,-3) {3};
  \node[lab] at (1,-0.35) {$i$};
  \node[lab] at (2,-0.35) {$i+1$};
  \node[lab] at (3,-0.35) {$i+2$};
  \draw[gridcol] (3.9,-1) grid (5.9,-3);
  \fill[black] (3.9,-1) circle (0.17);
  \filldraw[fill=white,draw=keptline,line width=0.6pt] (3.9,-2) circle (0.17);
  \filldraw[fill=white,draw=keptline,line width=0.6pt] (3.9,-3) circle (0.17);
  \fill[black] (4.9,-1) circle (0.17);
  \fill[black] (4.9,-2) circle (0.17);
  \filldraw[fill=white,draw=keptline,line width=0.6pt] (4.9,-3) circle (0.17);
  \filldraw[fill=white,draw=keptline,line width=0.6pt] (5.9,-1) circle (0.17);
  \fill[black] (5.9,-2) circle (0.17);
  \filldraw[fill=white,draw=keptline,line width=0.6pt] (5.9,-3) circle (0.17);
  \draw[draw=rectA,line width=1.2pt,line join=round] (4.9,-1) -- (5.9,-2) -- (4.9,-3) -- (3.9,-2) -- cycle;
  \node[capt,anchor=north] at (4.9,-3.4) {$aAb$};
  \draw[gridcol] (6.8,-1) grid (8.8,-3);
  \fill[black] (6.8,-1) circle (0.17);
  \filldraw[fill=white,draw=keptline,line width=0.6pt] (6.8,-2) circle (0.17);
  \filldraw[fill=white,draw=keptline,line width=0.6pt] (6.8,-3) circle (0.17);
  \fill[black] (7.8,-1) circle (0.17);
  \fill[black] (7.8,-2) circle (0.17);
  \filldraw[fill=white,draw=keptline,line width=0.6pt] (7.8,-3) circle (0.17);
  \filldraw[fill=white,draw=keptline,line width=0.6pt] (8.8,-1) circle (0.17);
  \filldraw[fill=white,draw=keptline,line width=0.6pt] (8.8,-2) circle (0.17);
  \fill[black] (8.8,-3) circle (0.17);
  \draw[draw=rectA,line width=1.2pt,line join=round] (7.8,-1) -- (8.8,-2) -- (7.8,-3) -- (6.8,-2) -- cycle;
  \node[capt,anchor=north] at (7.8,-3.4) {$aAc$};
  \draw[gridcol] (9.7,-1) grid (11.7,-3);
  \filldraw[fill=white,draw=keptline,line width=0.6pt] (9.7,-1) circle (0.17);
  \fill[black] (9.7,-2) circle (0.17);
  \filldraw[fill=white,draw=keptline,line width=0.6pt] (9.7,-3) circle (0.17);
  \fill[black] (10.7,-1) circle (0.17);
  \fill[black] (10.7,-2) circle (0.17);
  \filldraw[fill=white,draw=keptline,line width=0.6pt] (10.7,-3) circle (0.17);
  \fill[black] (11.7,-1) circle (0.17);
  \filldraw[fill=white,draw=keptline,line width=0.6pt] (11.7,-2) circle (0.17);
  \filldraw[fill=white,draw=keptline,line width=0.6pt] (11.7,-3) circle (0.17);
  \draw[draw=rectA,line width=1.2pt,line join=round] (10.7,-1) -- (11.7,-2) -- (10.7,-3) -- (9.7,-2) -- cycle;
  \node[capt,anchor=north] at (10.7,-3.4) {$bAa$};
  \draw[gridcol] (12.6,-1) grid (14.6,-3);
  \filldraw[fill=white,draw=keptline,line width=0.6pt] (12.6,-1) circle (0.17);
  \fill[black] (12.6,-2) circle (0.17);
  \filldraw[fill=white,draw=keptline,line width=0.6pt] (12.6,-3) circle (0.17);
  \fill[black] (13.6,-1) circle (0.17);
  \fill[black] (13.6,-2) circle (0.17);
  \filldraw[fill=white,draw=keptline,line width=0.6pt] (13.6,-3) circle (0.17);
  \filldraw[fill=white,draw=keptline,line width=0.6pt] (14.6,-1) circle (0.17);
  \fill[black] (14.6,-2) circle (0.17);
  \filldraw[fill=white,draw=keptline,line width=0.6pt] (14.6,-3) circle (0.17);
  \draw[draw=rectB,line width=1.2pt,line join=round,dash pattern=on 3pt off 2pt] (13.6,-1) -- (14.6,-2) -- (13.6,-3) -- (12.6,-2) -- cycle;
  \node[capt,anchor=north] at (13.6,-3.4) {$bAb$};
  \draw[gridcol] (15.5,-1) grid (17.5,-3);
  \filldraw[fill=white,draw=keptline,line width=0.6pt] (15.5,-1) circle (0.17);
  \fill[black] (15.5,-2) circle (0.17);
  \filldraw[fill=white,draw=keptline,line width=0.6pt] (15.5,-3) circle (0.17);
  \fill[black] (16.5,-1) circle (0.17);
  \fill[black] (16.5,-2) circle (0.17);
  \filldraw[fill=white,draw=keptline,line width=0.6pt] (16.5,-3) circle (0.17);
  \filldraw[fill=white,draw=keptline,line width=0.6pt] (17.5,-1) circle (0.17);
  \filldraw[fill=white,draw=keptline,line width=0.6pt] (17.5,-2) circle (0.17);
  \fill[black] (17.5,-3) circle (0.17);
  \draw[draw=rectB,line width=1.2pt,line join=round,dash pattern=on 3pt off 2pt] (16.5,-1) -- (17.5,-2) -- (16.5,-3) -- (15.5,-2) -- cycle;
  \node[capt,anchor=north] at (16.5,-3.4) {$bAc$};
  \draw[gridcol] (18.4,-1) grid (20.4,-3);
  \filldraw[fill=white,draw=keptline,line width=0.6pt] (18.4,-1) circle (0.17);
  \filldraw[fill=white,draw=keptline,line width=0.6pt] (18.4,-2) circle (0.17);
  \fill[black] (18.4,-3) circle (0.17);
  \fill[black] (19.4,-1) circle (0.17);
  \fill[black] (19.4,-2) circle (0.17);
  \filldraw[fill=white,draw=keptline,line width=0.6pt] (19.4,-3) circle (0.17);
  \fill[black] (20.4,-1) circle (0.17);
  \filldraw[fill=white,draw=keptline,line width=0.6pt] (20.4,-2) circle (0.17);
  \filldraw[fill=white,draw=keptline,line width=0.6pt] (20.4,-3) circle (0.17);
  \draw[draw=rectA,line width=1.2pt,line join=round] (19.4,-1) -- (20.4,-2) -- (19.4,-3) -- (18.4,-2) -- cycle;
  \node[capt,anchor=north] at (19.4,-3.4) {$cAa$};
  \draw[gridcol] (21.3,-1) grid (23.3,-3);
  \filldraw[fill=white,draw=keptline,line width=0.6pt] (21.3,-1) circle (0.17);
  \filldraw[fill=white,draw=keptline,line width=0.6pt] (21.3,-2) circle (0.17);
  \fill[black] (21.3,-3) circle (0.17);
  \fill[black] (22.3,-1) circle (0.17);
  \fill[black] (22.3,-2) circle (0.17);
  \filldraw[fill=white,draw=keptline,line width=0.6pt] (22.3,-3) circle (0.17);
  \filldraw[fill=white,draw=keptline,line width=0.6pt] (23.3,-1) circle (0.17);
  \fill[black] (23.3,-2) circle (0.17);
  \filldraw[fill=white,draw=keptline,line width=0.6pt] (23.3,-3) circle (0.17);
  \draw[draw=rectB,line width=1.2pt,line join=round,dash pattern=on 3pt off 2pt] (22.3,-1) -- (23.3,-2) -- (22.3,-3) -- (21.3,-2) -- cycle;
  \node[capt,anchor=north] at (22.3,-3.4) {$cAb$};
  \draw[gridcol] (24.2,-1) grid (26.2,-3);
  \filldraw[fill=white,draw=keptline,line width=0.6pt] (24.2,-1) circle (0.17);
  \filldraw[fill=white,draw=keptline,line width=0.6pt] (24.2,-2) circle (0.17);
  \fill[black] (24.2,-3) circle (0.17);
  \fill[black] (25.2,-1) circle (0.17);
  \fill[black] (25.2,-2) circle (0.17);
  \filldraw[fill=white,draw=keptline,line width=0.6pt] (25.2,-3) circle (0.17);
  \filldraw[fill=white,draw=keptline,line width=0.6pt] (26.2,-1) circle (0.17);
  \filldraw[fill=white,draw=keptline,line width=0.6pt] (26.2,-2) circle (0.17);
  \fill[black] (26.2,-3) circle (0.17);
  \draw[draw=rectA,line width=1.2pt,line join=round] (25.2,-1) -- (26.2,-2) -- (25.2,-3) -- (24.2,-2) -- cycle;
  \node[capt,anchor=north] at (25.2,-3.4) {$cAc$};
\end{tikzpicture}
}
\caption{The nine windows $sAt$ around a double $A$ in a word of class $\{A,C\}$, drawn on the
columns $i$, $i+1$, $i+2$ of the diamond $D(i,1)$. Black dots are blacked out and open dots are
white. The diamond is solid where the window passes and dashed where it fails.}
\label{fig:sixofnine}
\end{figure}

In a word of class $\{A,C\}$ we call a double \emph{isolated} if its column is not $1$ or $m$
and the other double is at least three columns away.

\begin{proposition}[separated doubles]\label{prop:dominant}
Let $m\ge6$. The number of valid words of class $\{A,C\}$ in which both doubles are isolated is
exactly $2\binom{m-4}26^23^{m-6}=8\binom{m-4}23^{m-4}$.
\end{proposition}

\begin{proof}
The doubles sit in columns $2\le p<q\le m-1$ with $q-p\ge3$, and $(p,q)\mapsto(p-1,q-3)$ is a
bijection onto the pairs $1\le p'<q'\le m-4$. So there are $\binom{m-4}2$ positions and two
orders of $A$ and $C$. Fix both. By Lemmas~\ref{lem:dichotomy}(a) and~\ref{lem:sixofnine} the
only conditions are on the two windows centered at the doubles. Each passes for six of the
nine choices of its two neighbors, and the four neighbors are different columns. The other
$m-6$ singletons are free, which gives $6^23^{m-6}$ valid words.
\end{proof}

So the class $\{A,C\}$ holds the fraction $(m-1)/(3m)$ of the $m^23^{m-1}$ candidates, and
all but $O(m)$ of its $\binom m2$ positions survive at the rate $(6/9)^2=4/9$. The product
$\frac13\cdot\frac49=\frac4{27}$ is the constant of Corollary~\ref{cor:threeasymp}.

\subsection{The dominant class and the exact formula}

Let $AC_m$ be the set of words $UAVCW$ of width $m$, with $U,V,W$ strings of singletons,
whose every window passes the test of Lemma~\ref{lem:diamond} with $F=\{A,C\}$, and put
$G_{AC}(z)=\sum_{m\ge2}|AC_m|z^m$. For $m\ge3$ these words are exactly the valid words of the
ordered class $AC$. We define $G_{AB}$ and $G_X$ in the same way from the words $UAVBW$ and
$UXW$ with $F=\{A,B\}$ and $F=\{A,B,C\}$. Put $R=1-3z$. A string of $n\ge1$ free singletons
has $3^n$ choices, and $3^{n-1}$ choices when its letter nearest to a double is prescribed,
so the prescribed strings have the generating function $\sum_{n\ge1}3^{n-1}z^n=z/R$.

\begin{proposition}[the class $\{A,C\}$]\label{prop:threeGAC}
The generating function of the ordered class $AC$ is
\begin{align}
G_{AC}(z)&=z^2\Bigl(1+\frac zR\Bigr)^2+z^3\Bigl(3+\frac{12z}R+\frac{13z^2}{R^2}\Bigr)
+\frac{9z^4(1-z)^2}{R^3}\label{eq:GAC}\\
&=\frac{z^2-4z^3+10z^4-8z^5-3z^6}{(1-3z)^3}.\nonumber
\end{align}
Its coefficients are $|AC_2|=1$, $|AC_3|=5$ and
\begin{equation}\label{eq:GACcoef}
|AC_m|=2(m-1)(3m+2)3^{m-5}\qquad(m\ge4).
\end{equation}
\end{proposition}

\begin{proof}
\emph{Step 1. The windows with both doubles.} In a window $tAC$ the diamond has $(i+1,1)$ and
$(i+2,2)$ black and $(i+1,3)$ white, and the tip $(i,2)$ is black exactly when $t=b$. For
$t=b$ it shows only $(i+1,3)$, and row (f) fails it through $2\in B_{i+1}$ and $C\in F$. For
$t\ne b$ it shows the oblique pair $(i,2)$, $(i+1,3)$, and row (e) asks whether $(i,3)$ and
$(i+1,2)$ are black. The second is, and the first is exactly when $t=c$. So $tAC$ passes only
for $t=a$. The reflection in the middle row followed by the reversal of the columns fixes the
ordered class $AC$ and swaps $a$ with $c$, so $ACt$ passes only for $t=c$. A window $AtC$ has
both tips black. For $t=b$ it shows the vertical pair, and for $t$ outer one point of column
$i+1$, and rows (c) and (f) pass both since $B\notin F$.

\emph{Step 2. The doubles are adjacent.} The word is $UACW$. The windows inside $U$ and $W$
and those with a double at a tip pass (Lemmas~\ref{lem:dichotomy}(a)
and~\ref{lem:sixofnine}(c)). By Step~1 the only conditions are that $U$, if nonempty, ends in
$a$, and that $W$, if nonempty, begins with $c$. So this case contributes $z^2(1+z/R)^2$.

\emph{Step 3. One singleton between the doubles.} The word is $UAtCW$, and the window $AtC$
passes. The window $uAt$, with $u$ the last letter of $U$, is governed by
Lemma~\ref{lem:sixofnine}(a). For $t=a$ every $u$ is allowed, for $t=b$ only $u=a$, and for
$t=c$ only $u\ne b$. So the numbers $p_t$ of allowed last letters of $U$ are $(3,1,2)$ for
$t=a,b,c$, and by Lemma~\ref{lem:sixofnine}(b) the numbers $q_t$ of allowed first letters of
$W$ are $(2,1,3)$. All other windows pass. The case contributes
\[
z^3\sum_{t}\Bigl(1+\frac{p_tz}R\Bigr)\Bigl(1+\frac{q_tz}R\Bigr)
=z^3\Bigl(3+\frac{12z}R+\frac{13z^2}{R^2}\Bigr),
\]
since $\sum p_t=\sum q_t=6$ and $\sum p_tq_t=6+1+6=13$.

\emph{Step 4. At least two singletons between the doubles.} The word is $UAVCW$ with
$V=v_1\cdots v_k$ and $k\ge2$. The windows that can fail are $uAv_1$ and $v_kCw$, with $u$
the last letter of $U$ and $w$ the first letter of $W$. If $U$ is empty, $v_1$ has three
choices. Otherwise the pair $(u,v_1)$ has nine choices, of which six are allowed. So the left
end contributes $3+6z/R$, and the right end likewise. The letters $v_2,\dots,v_{k-1}$ are
free, with $\sum_{k\ge2}3^{k-2}z^{k-2}=1/R$. So the case contributes
\[
z^4\Bigl(3+\frac{6z}R\Bigr)^2\frac1R=\frac{9z^4(R+2z)^2}{R^3}=\frac{9z^4(1-z)^2}{R^3}.
\]

\emph{Step 5. The coefficients.} The three cases are disjoint and exhaustive, which gives the
first form in~\eqref{eq:GAC}, and putting the terms over $R^3$ gives the second. Since
$[z^n]R^{-3}=\binom{n+2}23^n$, the coefficient of $z^m$ is
\begin{multline*}
\binom m23^{m-2}-4\binom{m-1}23^{m-3}+10\binom{m-2}23^{m-4}\\
-8\binom{m-3}23^{m-5}-3\binom{m-4}23^{m-6},
\end{multline*}
where a binomial coefficient with a negative top entry is $0$. For $m\ge4$ every top entry is
nonnegative, and the expression is the polynomial identity $2(m-1)(3m+2)3^{m-5}$, which
is~\eqref{eq:GACcoef}. For $m=2$ and $m=3$ some top entries are negative, and the values are
$1$ and $5$.
\end{proof}

The other classes are exponentially smaller, because $B\in F$ forbids the singleton windows
$bab$, $bbb$ and $bcb$ (Lemma~\ref{lem:dichotomy}(b)). A string of singletons with no two
letters $b$ at distance two is a \emph{restricted run}.

\begin{lemma}[runs]\label{lem:runs}
Let $t_n$ be the number of words of length $n$ in the letters $a,b,c$ with no two consecutive
letters $b$. Then $t_0=1$, $t_1=3$ and $t_n=2t_{n-1}+2t_{n-2}$ for $n\ge2$. Put
$\rho=1+\sqrt3$ and $\kappa=9/\rho^2$. Then $t_n\le(3/\rho)\rho^n$ for every $n\ge0$, and the
number of restricted runs of length $n$ is at most $\kappa\rho^n$.
\end{lemma}

\begin{proof}
A word counted by $t_n$ ends in $a$ or $c$, with $2t_{n-1}$ choices, or in $b$ preceded by $a$
or $c$, with $2t_{n-2}$ choices. The bound holds at $n=0$ and $n=1$, and since
$\rho^2=2\rho+2$ the induction step is $t_n\le(3/\rho)(2\rho^{n-1}+2\rho^{n-2})=(3/\rho)\rho^n$.
In a restricted run the letters in the odd positions form a word with no two consecutive
letters $b$, and so do the letters in the even positions. So a restricted run of length $n$
has at most $t_{\lceil n/2\rceil}t_{\lfloor n/2\rfloor}\le(3/\rho)^2\rho^n=\kappa\rho^n$
choices.
\end{proof}

\begin{theorem}[the number of three-row optima]\label{thm:threecount}
For $m\ge3$ let $E_m$ be the number of maximum blackouts of the $3\times m$ grid whose word
lies outside the ordered classes $AC$ and $CA$. So $E_m=4|AB_m|+|X_m|$ counts the four
ordered classes $AB$, $BA$, $BC$, $CB$ and the triple class.
\begin{enumerate}[(a)]
\item For every $m\ge4$,
\begin{equation}\label{eq:threecount}
a_3(m)=4(m-1)(3m+2)3^{m-5}+E_m.
\end{equation}
\item For every $m\ge3$, with $\rho=1+\sqrt3$ and $\kappa=9/\rho^2$,
\begin{equation}\label{eq:Embound}
0\le E_m\le4\binom m2\kappa^3\rho^{m-2}+m\kappa^2\rho^{m-1}.
\end{equation}
\item At $m=3$ the classes contribute $a_3(3)=53=2\cdot5+4\cdot7+15$.
\end{enumerate}
\end{theorem}

\begin{proof}
(a) The reversal of the columns and the reflection in the middle row map rectangles to
rectangles and presentations to presentations, so they map maximum blackouts to maximum
blackouts. The reversal reverses the word, so it maps the ordered class $AC$ onto $CA$, $AB$
onto $BA$ and $BC$ onto $CB$. The reflection swaps $a$ with $c$ and $A$ with $C$ and fixes
$B$ and $X$, so it maps $AB$ onto $CB$. Hence $|AC_m|=|CA_m|$ and
$|AB_m|=|BA_m|=|BC_m|=|CB_m|$, and the seven classes give
\begin{equation}\label{eq:classsum}
a_3(m)=2|AC_m|+4|AB_m|+|X_m|=2|AC_m|+E_m\qquad(m\ge3).
\end{equation}
Now~\eqref{eq:GACcoef} gives~\eqref{eq:threecount} for $m\ge4$.

(b) The lower bound holds since $E_m$ is a count. For the upper bound we drop every condition
except the three forbidden singleton windows of Lemma~\ref{lem:dichotomy}(b), which hold in
every class counted by $E_m$ because $B\in F$. Take a word of one of the four ordered double
classes. It is determined by the positions of its two doubles, with $\binom m2$ choices, and
by its three singleton strings. These are restricted runs of total length $m-2$, with at most
$\kappa^3\rho^{m-2}$ choices together (Lemma~\ref{lem:runs}). A word of the triple class is
determined by the position of $X$, with $m$ choices, and by two restricted runs of total length
$m-1$, with at most $\kappa^2\rho^{m-1}$ choices.

(c) The class counts at width $3$ are $|AC_3|=5$ from Proposition~\ref{prop:threeGAC} and
$|AB_3|=7$, $|X_3|=15$ from the coefficients of $z^3$ in Propositions~\ref{prop:GAB}
and~\ref{prop:GX}.
\end{proof}

The formula~\eqref{eq:threecount} is exact for $m\ge4$, and the remainder $E_m$ is a count,
not a signed error. At $m=3$ the right side of~\eqref{eq:GACcoef} would be $44/9$, which is
why the theorem starts at $m=4$. Theorem~\ref{thm:threeclosed} gives a closed form of $E_m$.

\subsection{The generating function}

Put
\[
R=1-3z,\qquad L=1-2z-2z^2,\qquad N=1+2z^2,\qquad D=LN.
\]

\begin{theorem}[three-row generating function, parts (c) and (d) computer-assisted]\label{thm:threegf}
Let $Q=R^3L^3N^3$.
\begin{enumerate}[(a)]
\item For every $m\ge3$,
\begin{gather*}
a_3(m)=2|AC_m|+4|AB_m|+|X_m|,\\
\sum_{m\ge3}a_3(m)z^m=2G_{AC}+4G_{AB}+G_X-z-12z^2.
\end{gather*}
\item The function $G_{AC}$ is~\eqref{eq:GAC}. The functions $G_{AB}$ and $G_X$ are rational
with denominators dividing $D^3$ and $D^2$, and Propositions~\ref{prop:GAB} and~\ref{prop:GX}
give them exactly.
\item The reduced denominator of $\sum_{m\ge3}a_3(m)z^m$ is exactly
\[
Q(z)=(1-3z)^3(1-2z-2z^2)^3(1+2z^2)^3.
\]
The exact orders at the roots of $L$ and $N$ rest on the computer-assisted polynomial
form~\eqref{eq:PB} of $G_{AB}$.
\item Put
\begin{align*}
P(z)={}&z-3z^2-34z^3+228z^4-641z^5+1307z^6-2186z^7\\
&+2056z^8-236z^9-1060z^{10}+2584z^{11}-4264z^{12}\\
&-1440z^{13}+704z^{14}-1024z^{15}+2080z^{16}+2176z^{17}+384z^{18}.
\end{align*}
Then
\begin{equation}\label{eq:PQ}
\sum_{m\ge3}a_3(m)z^m=\frac{P(z)}{Q(z)}-z-12z^2.
\end{equation}
\end{enumerate}
\end{theorem}

\begin{proof}
(a) The first identity is~\eqref{eq:classsum}. The three generating functions count all words
of their forms, and for $m\ge3$ these are exactly the valid words of the ordered classes, so
the series agree from $z^3$ on. The coefficients of $z$ and $z^2$ in $2G_{AC}+4G_{AB}+G_X$ are
$1$ and $12$. They count the word $X$, the six words $sX$ and $Xs$ with $s$ a singleton, and
the six words made of two different doubles. The terms $z+12z^2$ remove them.

(b) is Proposition~\ref{prop:threeGAC} with Propositions~\ref{prop:GAB} and~\ref{prop:GX} of
Appendix~\ref{app:threegf}, and (c) is Theorem~\ref{thm:denominator} there. In short,
$G_{AB}$ has a pole of order three at every root of $L$ and of $N$, which the other two
classes cannot cancel. The function $G_{AC}$ has one at $z=1/3$, where the other two classes
are analytic.

(d) Substituting the formulas of (b) into (a) and clearing denominators is exact polynomial
algebra, and it is computer-assisted. The script \texttt{three\_row\_closed\_form.py} does it
with SymPy and obtains the displayed $P$, with $P$ and $Q$ coprime. The script
\texttt{three\_row\_closed\_form\_check.py} checks~\eqref{eq:PQ} against the transfer matrix of
Remark~\ref{rem:transfer}, in exact integer arithmetic, through width $100$.
\end{proof}

\begin{corollary}[asymptotics]\label{cor:threeasymp}
With $\rho=1+\sqrt3$, as $m\to\infty$,
\begin{align}
a_3(m)&=\frac4{81}m^23^m-\frac4{243}(m+2)3^m+O(m^2\rho^m),\label{eq:asymp1}\\
\frac{a_3(m)}{m^23^{m-1}}&=\frac4{27}-\frac4{81m}-\frac8{81m^2}+O\bigl((\rho/3)^m\bigr).
\label{eq:asymp2}
\end{align}
In particular $a_3(m)\sim\frac4{81}m^23^m$, and the fraction of the $m^23^{m-1}$ candidates of
Proposition~\ref{prop:threeclass} that are valid tends to $4/27$.
\end{corollary}

\begin{proof}
The main term of~\eqref{eq:threecount} is
$4(3m^2-m-2)3^{m-5}=\frac4{81}m^23^m-\frac4{243}(m+2)3^m$, and~\eqref{eq:Embound} gives
$E_m=O(m^2\rho^m)$, which is~\eqref{eq:asymp1}. Dividing by $m^23^{m-1}$
gives~\eqref{eq:asymp2}, since $\frac4{243}(m+2)3^m/(m^23^{m-1})=\frac4{81m}+\frac8{81m^2}$.
\end{proof}

The second route to the asymptotic is the pole of order three of~\eqref{eq:PQ} at $z=1/3$,
where $(1-3z)^3P/Q$ tends to $8/81$, and the two routes agree. The error term
in~\eqref{eq:asymp1} has the right order, and the partial fractions of $4G_{AB}+G_X$ give it
exactly.

\begin{theorem}[three-row closed form, computer-assisted exact algebra]\label{thm:threeclosed}
Let $\rho=1+\sqrt3$ and $\sigma=1-\sqrt3$. For every $m\ge3$,
\begin{equation}\label{eq:Eclosed}
E_m=p_+(m)\rho^m+p_-(m)\sigma^m+T_m,
\end{equation}
where
\[
p_\pm(m)=\frac{(19\pm11\sqrt3)m^2+(61\pm41\sqrt3)m+(104\pm28\sqrt3)}{432}
\]
and
\[
T_m=\begin{cases}
(-1)^{m/2}\,2^{m/2}\,\dfrac{2m^2+41m-104}{216}&\text{if $m$ is even,}\\[2ex]
(-1)^{(m-1)/2}\,2^{(m-1)/2}\,\dfrac{7m^2-32m-175}{432}&\text{if $m$ is odd.}
\end{cases}
\]
Hence $E_m\sim\frac{19+11\sqrt3}{432}m^2\rho^m$, and $E_m>0$ for every $m\ge4$.
\end{theorem}

\begin{proof}
By Theorem~\ref{thm:threegf}(a) and~(b), $4G_{AB}+G_X$ is a rational function with the
denominator $L^3N^3$, and its coefficient of $z^m$ is $E_m$ for $m\ge3$. Its partial-fraction
expansion has a polynomial part of degree one, and the rest is~\eqref{eq:Eclosed}. The roots
of $L^3N^3$ are $1/\rho$, $1/\sigma$ and $\pm1/(i\sqrt2)$, each of multiplicity three, so the
coefficients are quadratic polynomials in $m$, and $T_m$ collects the two conjugate terms at
$\pm i\sqrt2$. This step is computer-assisted. The script
\texttt{three\_row\_closed\_form.py} solves for the twelve rational coefficients and verifies
the identity by clearing denominators symbolically. Without a computer algebra system, the
script \texttt{three\_row\_closed\_form\_check.py} evaluates the right side in exact rational
arithmetic over $\mathbb Q(\sqrt3)$ and compares it with the transfer counts of
Remark~\ref{rem:transfer} for every $2\le m\le100$. Since both sides satisfy the order-$12$
recurrence of $L^3N^3$ from $m=14$ on, the agreement at $m=2,\dots,13$ already proves the
identity. For the positivity, $p_+(m)\ge19m^2/432$, while $|p_-(m)|\le m^2$, $|\sigma|<1$ and
$|T_m|\le2^{m/2}m^2$ for $m\ge4$. So $E_m\ge m^2(19\rho^m/432-1-2^{m/2})$, and the bracket is
positive for $m\ge6$, since $19\rho^6/432>18>1+2^3$ and $\rho>\sqrt2$. The values $E_4=168$
and $E_5=620$ are positive.
\end{proof}

\begin{corollary}[a recurrence of order $15$, computer-assisted]\label{cor:threerec}
Write $Q(z)=\sum_{j=0}^{15}q_jz^j$, so that
\begin{align*}
(q_0,\dots,q_7)&=(1,-15,93,-317,702,-1188,1512,-972),\\
(q_8,\dots,q_{15})&=(-144,1136,-2640,1104,800,576,3456,1728).
\end{align*}
Then
\begin{equation}\label{eq:rec}
\sum_{j=0}^{15}q_ja_3(m-j)=0\qquad\text{for every }m\ge19,
\end{equation}
and $19$ is the exact threshold, since the left side equals $384$ at $m=18$. The values at
widths $3$ to $18$ are $53$, $224$, $892$, $3420$, $12704$, $45864$, $161992$, $562912$,
$1930312$, $6545496$, $21988784$, $73296272$, $242715784$, $799199592$, $2618729664$,
$8544429216$.
\end{corollary}

\begin{proof}
By~\eqref{eq:PQ}, $Q(z)\sum_{m\ge3}a_3(m)z^m=P(z)-(z+12z^2)Q(z)$, a polynomial of degree $18$
with the leading coefficient $384$ of $P$, since $(z+12z^2)Q$ has degree $17$. Comparing the
coefficients of $z^m$ gives~\eqref{eq:rec} for $m\ge19$ and the value $384$ at $m=18$. The
values at widths $3$ to $18$ are the coefficients of~\eqref{eq:PQ}.
\end{proof}

The two programs of Section~\ref{sec:verify} give the values at widths $3$ to $9$. The
integer program finds them, and the exhaustive check from the definition confirms them. The
values at widths $10$ to $18$ follow from the generating function, and we have not enumerated
them from the definition. The counts of each class at widths $3$ to $10$, from a test of
every candidate against every rectangle, are in the data file
\texttt{three\_row\_definition\_check.json}. They agree with the series of
Propositions~\ref{prop:threeGAC}, \ref{prop:GAB} and~\ref{prop:GX}, and no proof uses them.

\begin{remark}[the $350$-state transfer, computer-assisted]\label{rem:transfer}
Before the class decomposition we obtained~\eqref{eq:PQ} by a transfer matrix, and it is an
independent confirmation. A state records a target pair set $F$ among the four possible ones,
the set $U\subseteq F$ of pairs seen so far, the excess $\sum(|B_i|-1)\in\{0,1,2\}$, and the
last two letters. A letter is rejected in four cases. It brings a pair outside $F$, repeats a
pair of $U$, raises the excess above $2$, or completes a window that fails
Lemma~\ref{lem:diamond} with the target $F$. A word is accepted when $U=F$ and the excess is
$2$, under its own pair set only. There are $350$ reachable states, $85$ of them accepting,
and $1134$ transitions. With $M$ the transition matrix, $\alpha$ the four start states and
$\beta$ the accepting ones, $a_3(m)=\alpha M^m\beta$ for $m\ge3$. The script checks, in exact
integer arithmetic, that $\sum_{j=0}^{15}q_jM^{19-j}\beta=0$ in all $350$ coordinates, which
proves~\eqref{eq:rec} for $m\ge19$ directly, and that $\alpha(I-zM)^{-1}\beta$
agrees with~\eqref{eq:PQ} through $z^{100}$, class by class.
\end{remark}

\subsection{Optima up to symmetry}

The counts $a_3(m)$ are labeled. The symmetries of the $3\times m$ grid form the Klein
four-group generated by the reflection in the middle row and the reversal of the columns,
whose product is the half turn. On words the reflection swaps $a$ with $c$ and $A$ with $C$,
the reversal reverses the word, and the half turn does both. Let $o_3(m)$ be the number of
orbits of the maximum blackouts under this group, and let $t_n$ be as in Lemma~\ref{lem:runs}.

\begin{proposition}[optima up to symmetry, values computer-assisted]\label{prop:threeorbits}
Let $m\ge3$ and let $r_m$, $v_m$ and $w_m$ be the numbers of maximum blackouts of the
$3\times m$ grid fixed by the reflection, the reversal and the half turn.
\begin{enumerate}[(a)]
\item $o_3(m)=\bigl(a_3(m)+r_m+v_m+w_m\bigr)/4$.
\item $r_m=1$ if $m=3$ and $r_m=0$ otherwise. The only reflection-fixed optimum is $bXb$.
\item $v_m=0$ for even $m$, $v_3=1$, $v_5=2$ and, for $m=2k+1$ with $k\ge3$,
$v_m=4\,t_{\lceil k/2\rceil-2}\,t_{\lfloor k/2\rfloor-1}$. The reversal-fixed optima lie in
the triple class.
\item $w_3=5$ and, for $m\ge4$ with $n=\lfloor m/2\rfloor$,
$w_m=4n3^{n-2}+x_m$, where $x_m=0$ for even $m$ and
$x_m=2\,t_{\lceil k/2\rceil}\,t_{\lfloor k/2\rfloor-1}$ for $m=2k+1$ with $k\ge2$. The
half-turn-fixed optima lie in the class $\{A,C\}$ and in the triple class.
\item The values $o_3(3),\dots,o_3(15)$ are $15$, $58$, $227$, $864$, $3190$, $11502$,
$40549$, $140863$, $482755$, $1636860$, $5497794$, $18325769$, $60680951$.
\end{enumerate}
\end{proposition}

So $o_3(m)=a_3(m)/4+O(m3^{m/2})$, and almost every optimum has a trivial stabilizer. At
$m=3$ the orbit of $bXb$ has one element, the orbits of $AbC$ and $aXc$ have two, and the
other twelve orbits have four. At $m=4$ the eight half-turn-fixed words are $AacC$, $AbbC$,
$AcaC$, $CacA$, $CbbA$, $CcaA$, $aACc$ and $cCAa$. They form four orbits of size two, and the
other $54$ orbits have size four.

\begin{proof}
(a) is Burnside's lemma, since the symmetries map maximum blackouts to maximum blackouts. The
fixed words in (b) to (d) follow from the forbidden windows of Lemmas~\ref{lem:dichotomy},
\ref{lem:sixofnine}, \ref{lem:threewindows} and Step~1 of the proof of
Proposition~\ref{prop:threeGAC}. The values in (e) rest on Corollary~\ref{cor:threerec}. For
$m\le10$ the script \texttt{three\_row\_orbits.py} also counts the orbits directly on the
optima found from the definition, and the counts agree.

(b) The letters fixed by the reflection are $b$, $B$ and $X$, and two doubles would both be
$B$. So a fixed candidate is $b^iXb^j$, and the forbidden windows $bbX$ and $Xbb$ force
$i,j\le1$ and $m=3$.

(c) A palindrome cannot carry two different doubles, since the double at column $i$ would
equal the one at column $m+1-i$. So $m=2k+1$ and the word is $wXw^R$ with $w$ of length $k$.
The windows $aXa$ and $cXc$ are forbidden and $bXb$ is allowed, so $w$ ends in $b$, and by the
window $bsX$ its second-to-last letter is outer. Otherwise $w$ is an arbitrary restricted run.
Its letters in the positions $k,k-2,\dots$ form a word with no two consecutive letters $b$
that ends in $b$, with $2t_{\lceil k/2\rceil-2}$ choices for $k\ge3$. Its letters in the
positions $k-1,k-3,\dots$ form one that ends in an outer letter, with
$2t_{\lfloor k/2\rfloor-1}$ choices. The cases $k=1,2$ give $b$ and $ab$, $cb$.

(d) The half turn sends the letter at column $i$ to the reflected letter at column $m+1-i$. A
fixed word of the triple class is $wX\overline{w}^R$, with $\overline{w}$ the reflection of
$w$ and $m=2k+1$. The windows $aXc$, $bXb$ and $cXa$ are allowed, so the only condition is
that $w$ is a restricted run with an outer second-to-last letter, which gives $x_m$ by the
parity split of (c). A fixed word with doubles has them at columns $p$ and $m+1-p$, one $A$
and one $C$, so it is $UAVC\overline U^R$ with $V=\overline V^R$, or the same with $C$ first.
Take the order $AC$, let $u=|U|$ and $k=|V|$, so $m=2u+k+2$ and $n=u+\lfloor k/2\rfloor+1$.
By Step~1 of the proof of Proposition~\ref{prop:threeGAC} and Lemma~\ref{lem:sixofnine} the
conditions are these. If $k\le1$, then $V$ is empty or $V=b$, and the letter of $U$ next to
$A$ is $a$. If $k\ge2$, the pair of that letter and the first letter of $V$ avoids $bb$, $bc$
and $cb$. Every other window passes or is the half-turn image of one of these. The word $V$ has $\lfloor k/2\rfloor$ free
letters, since its middle letter, for odd $k$, is $b$. So the count is $3^{n-2}$ for $k\le1$,
$3^{n-1}$ for $u=0$ and $k\ge2$, and $2\cdot3^{n-2}$ for each of the $n-2$ pairs $(u,k)$ with
$u\ge1$ and $k\ge2$. The sum is $2n3^{n-2}$ for $m\ge4$, and the order $CA$ gives the same.
At $m=3$ the fixed words are $AbC$, $CbA$, $aXc$, $bXb$ and $cXa$.
\end{proof}
